\subsection{LOCALLY CLOSED SUBSPACES}

DEFINITION 2. A subset L of a topological space X is said to be locally closed at a point \(x \in L\) if there is a neighbourhood V of x in X such that \(L \cap V\) is a closed subset of the subspace V. L is said to be locally closed in X if it is locally closed at each \(x \in L\) .

Remark. Let F be a subset of X such that for each point x of X there is a neighbourhood V of x such that \(V \cap F\) is closed in the subspace V; then it follows from Proposition 3 of no. 1 that F is closed in X. On the other hand, Proposition 5 below shows that in general there are locally closed sets which are not closed in X.

PROPOSITION 5. For a subset L of a topological space X, the following properties are equivalent:

\begin{enumerate}
\def\labelenumi{\alph{enumi})}
\item
  L is locally closed in X.
\item
  L is the intersection of an open subset and a closed subset of X.
\item
  L is open in its closure \(\overline{L}\) in X.
\end{enumerate}

If L is locally closed, then, for each \(x \in L\) , there is an open neighbourhood \(V_{x}\) of x in X such that \(L \cap V_{x}\) is closed in \(V_{x}\) ; \(U = \bigcup_{x \in L} V_{x}\) is open in X, and Proposition 3 of no. 1 shows that L is closed in U; therefore a) implies b). If \(L = U \cap F\) , where U is open and F closed in X, we have \(\overline{L} \subset F\) ; hence \(L \subset U \cap \overline{L} \subset U \cap F = L\) , which shows that \(L = U \cap \overline{L}\) is open in \(\overline{L}\), so that b) implies c). Finally, if \(L = U \cap \overline{L}\) , where U is open in X, \(\overline{L}\) is closed in U, hence locally closed, and thus c) implies a).

COROLLARY. Let \(f: \mathbf{X} \to \mathbf{X}'\) be a continuous map and \(\mathbf{L}'\) a locally closed subset of \(\mathbf{X}'\) ; then \(\overline{f}^{1}(\mathbf{L}')\) is locally closed in \(\mathbf{X}\) .

This follows immediately from Proposition 5 above and Theorem 1 of § 2, no.1.

